\documentclass[../main.tex]{subfiles}

\begin{document}

Durante la mayor parte del texto podremos usar tipos en vez de conjuntos y
tratarlos intuitivamente de la misma manera. Sin embargo, llegará un momento,
al definir la semántica denotacional, en el que las diferencias entre conjuntos
y tipos será crucial.
En teoría de tipos se puede definir un conjunto como una función que toma un
elemento y devuelve una proposición, que indica si el elemento pertenece o no al
conjunto.

\begin{minted}{lean4}
def Set (α: Type u) := α → Prop

def setOf (p: α → Prop): Set α :=
  p
\end{minted}


Podemos ahora \enquote{demostrar} un axioma clásico para conjuntos, el 
axioma de extensionalidad, que dice que dos conjuntos son iguales si y solo si
tienen los mismos elementos. Esto lo hacemos gracias a dos axiomas,
el de extensionalidad funcional y el de extensionalidad proposicional.

\begin{minted}{lean4}
theorem ext {a b: Set α} (h: ∀(x: α), x ∈ a ↔ x ∈ b):
  a = b :=
  funext (λ x ↦ propext (h x))
\end{minted}



\end{document}